% Part: computability
% Chapter: recursive-functions
% Section: partial-functions

\documentclass[../../../include/open-logic-section]{subfiles}

\begin{document}

\olfileid{cmp}{rec}{par}
\olsection{部分递归函数}

作为递归函数定义的动机，请注意，我们关于存在非原始递归的可计算函数的证明实际上确立了更多的结论。
论证很简单：我们所利用的全部事实只是，可以枚举函数 $f_0,f_1,\dots$，使得将 $f_x(y)$ 视为 $x$ 和 $y$ 的函数时，它是可计算的。
因此，该论证适用于\emph{任何能以这种方式枚举的函数类}。
这使我们陷入一个困境：我们想要明确地描述可计算函数；但是，对可计算函数集合的任何明确描述都不可能是穷尽的！

解决之道是引入\emph{偏}函数。
我们将看到，偏可计算函数\emph{的确}是可枚举的。\iftag{TMs}{事实上，我们已大致知道情况确实如此，因为Turing机是可以系统地枚举的。}{}
我们稍后将回到对角论证，并探讨为何在包含偏函数时它就无法成立。

现在的问题是：我们需要在原始递归函数的基础上添加什么，才能得到所有的偏递归函数？我们需要做两件事：

\begin{enumerate}
\item 修改原始递归函数的定义，使其同样允许偏函数。
\item 向该定义中\emph{添加}一些内容，从而包含新的偏函数。
\end{enumerate}

第一点很容易。和之前一样，我们将从零函数、后继函数和投影函数出发，并在复合与原始递归下封闭。唯一的区别是，我们必须修改复合与原始递归的定义，以允许定义中的某些项可能无定义的情况。如果 $f$ 和 $g$ 是偏函数，我们将用 $f(x) \fdefined$ 表示 $f$ 在 $x$ 处有定义，即 $x$ 在 $f$ 的定义域中；而用 $f(x) \fundefined$ 表示相反的情况，即 $f$ 在 $x$ 处无定义。我们将用 $f(x) \simeq g(x)$ 表示要么 $f(x)$ 和 $g(x)$ 都无定义，要么它们都有定义并且相等。我们也将这些记号用于更复杂的项。我们将采用如下约定：如果 $h$ 和 $g_0$, \dots, $g_k$ 都是偏函数，那么
\[
h(g_0(\vec x),\dots,g_k(\vec x))
\]
有定义当且仅当每个 $g_i$ 都在 $\vec x$ 处有定义，并且 $h$ 在 $g_0(\vec x)$, \dots, $g_k(\vec x)$ 处有定义。基于这一理解，偏函数的复合与原始递归定义就与上文完全相同，只不过我们必须将“$=$”替换为“$\simeq$”。

为了得到偏函数，我们将向原始递归函数定义中添加的是\emph{无界搜索算子}。如果 $f(x,\vec z)$ 是自然数上的任意偏函数，定义 $\mu x \; f(x,\vec z)$ 为

\begin{quote}
  使得 $f(0,\vec z), f(1,\vec z), \dots, f(x,\vec
  z)$ 均有定义且 $f(x,\vec z) = 0$ 的最小 $x$（如果这样的 $x$ 存在），
\end{quote}

并约定否则 $\mu x \; f(x,\vec z)$ 无定义。这唯一地定义了 $\mu x \; f(x,\vec z)$。

\begin{explain}
请注意，我们的定义并未提及\iftag{TMs}{Turing机，或}{}算法，或任何特定的计算模型。但是，与复合和原始递归一样，无界搜索背后也有一种操作性和计算性的直观。就偏函数的可计算性而言，函数无定义的自变元对应于计算不会停机的输入。计算 $\mu x \;
f(x,\vec z)$ 的过程相当于：计算 $f(0,\vec z), f(1,\vec z),
f(2,\vec z)$，直到返回值 0。然而，如果任何中间计算不停机，那么 $\mu x \; f(x,\vec z)$ 的计算也不会停机。
\end{explain}

如果 $R(x,\vec z)$ 是任意关系，$\mu x \; R(x,\vec z)$ 被定义为 $\mu x \; (1 \tsub \Char{R}(x,\vec z))$。换言之，$\mu x \;
R(x,\vec z)$ 返回使 $R(x,\vec z)$ 成立的最小 $x$ 值。因此，如果 $f(x,\vec z)$ 是全函数，那么 $\mu x \; f(x,\vec
z)$ 就等同于 $\mu x \; (f(x,\vec z) = 0)$。但请注意，我们最初的定义更为一般，因为它允许 $f(x,\vec z)$ 可能并非处处有定义（相反，关系的特征函数总是全函数）。

\begin{defn}
\emph{部分递归函数}集合是满足以下条件的最小偏函数（从自然数到自然数，具有不同元数）集合：包含零函数、后继函数和投影函数，且在复合、原始递归和无界搜索下封闭。
\end{defn}

当然，部分递归函数中有一些恰好是全函数，即对每个自变元都有定义。

\begin{defn}
\ollabel{defn:recursive-fn}
\emph{递归函数}集合是所有作为全函数的部分递归函数构成的集合。
\end{defn}

递归函数有时也被称为“全递归函数”，以强调它是处处有定义的。

\end{document}